\documentclass[a4paper,12pt]{ctexart}
\usepackage{xcolor}
\usepackage{graphicx}
\usepackage[top=1.8cm, bottom=1.8cm, left=1.8cm, right=1.8cm]{geometry}
\usepackage{array}
\usepackage{multirow}
\usepackage{hyperref}
\hypersetup{colorlinks=true,linkcolor=blue!70!black}
\linespread{1.35}

\newcommand{\ruby}[2]{#1\textsuperscript{#2}}
\newcommand{\shihai}[2]{%
  \vspace{0.35cm}
  \noindent
  \fcolorbox{orange!60!black}{orange!8}{%
    \begin{minipage}{0.915\textwidth}
    \textbf{\textcolor{orange!70!black}{地理笔记：}#1} \\[0.1cm]
    {\small #2}
    \end{minipage}%
  }
  \vspace{0.35cm}
}

\title{\textcolor{blue!70!black}{\Huge\textbf{经济与人文地理}}}
\author{}
\date{}

\begin{document}
\maketitle
\thispagestyle{empty}

\tableofcontents
\newpage

\section{\textcolor{orange!80!black}{第一课\quad 读地图：方向、比例尺、图例和经纬度}}

\subsection*{一张地图比你想象的说了更多}

拿起你的手机打开地图应用。你能看到你家在哪个街区，到学校那条蓝色的步行路线拐了几道弯。但你有没有想过一个问题——这张照片一样的平面图是怎么把一颗球体的表面摊平的？

地球是一个近似椭球体。把一个球体完整地铺在一张纸上而不产生任何拉伸、撕裂或变形——从数学上讲\textbf{不可能}。每一张世界地图都在做一些妥协。墨卡托投影（就是你在大多数教室墙上看到的那种长方形世界地图）保留了方向，但代价是\textbf{把高纬度地区面积严重放大了}。在墨卡托投影里，格陵兰岛看起来几乎和非洲一样大——但实际格陵兰的面积约216万平方公里，非洲的面积约3037万平方公里，后者是前者的\textbf{十四倍}。你看到的那张地图在悄悄地欺骗你的眼睛。

\subsection*{三个读图的基本零件}

不管你看的是什么地图——一张地铁线路图、一张地形图、一张景区导览图——它一定带着三个基本组成元素：

\textbf{方向。}大部分现代地图默认上北下南、左西右东。这不是一条自然规律——澳大利亚人做地图时完全可以把南放在上面——但全球统一化之后绝大多数地图都接受了这条约定。在地面实际操作时，如果你找不到北，看太阳：北半球正午太阳在南方（你的影子指向北方）。如果你手里有一块手表：将时针指向太阳，时针和12点刻度之间的中间方向\u2014\u2014大约就是南。

\textbf{比例尺。}地图下方的那一条短线加数字——"1:50000"\u2014\u2014\:意思是图上的一厘米代表实际地面上的五万厘米，也就是500米。如果比例尺很大（分母很小，比如1:1000），你能看到你家门口的人行横道。如果比例尺很小（分母极大，比如1:5000000），一整座城市压缩成了脚底的一块斑。

\textbf{图例。}地图边上的一排小方块——蓝色线是河流、黑色双线是铁路、绿色块是公园、红色虚线是地铁——没有这些标识，地图就是一堆颜料的拼接。

\subsection*{等高线——怎样把一座山画在纸上}

地理图一个最重要的技能是\textbf{看懂地形}。等高线把\textbf{海拔高度相等的所有点连成一条闭环曲线}。线越密——坡度越陡。线稀疏——地面坡度平缓。一条线圈起来的范围如果比内部的圈儿小——说明那是一座隆起（山峰）。如果圈与圈之间向外展开并且每一圈的海拔往下掉——这是一个洼地。

在等高线地图上，你可以在心里重建出三维地形，判断大江大河必然从高处流向低处，判断公路该绕山谷盘不盘、哪一侧地势是屏障。

\shihai{经度和纬度——你家的球面坐标}{
  把地球想象成一个西瓜。从瓜蒂到瓜底切片的那些纵切线是\textbf{经线}（子午线）。所有的经线都在南北两极汇聚。0度经线——人为指定的起点——穿过英国伦敦郊区的格林尼治皇家天文台。北京的经度大约是东经116度\u2014\u2014"东"说明从格林尼治往东数116度。

  横着切的那些平行的圈是\textbf{纬线}。赤道是最长的纬线圈（纬度0度），往两侧走到南北极各自缩成点（纬度90度）。北京的纬度大约是北纬40度。随纬度变化最大的三件东西是：\textbf{温度}\u2014\u2014每往北推大约111公里（1度纬度），年平均气温大约下降0.6℃；\textbf{白昼长度和四季}{\footnotesize }。哈尔滨（北纬45度）的夏至白天比广州（北纬23度）长将近一个半小时。

  有了经纬度，你可以在球面上钉住任何一点。你家的那组两位数坐标——不是随机数，是用球面数学算出来的唯一地址。
}

\subsection*{要点总结}
\begin{itemize}
  \item 把球面地图摊平成平面一定会产生变形——墨卡托投影放大了高纬度地区（格陵兰看起来和非洲一样大，实际只有非洲的十四分之一）
  \item 读地图的三个基本要素：方向（上北下南是约定）、比例尺（图上距离与实际距离的比值）、图例（符号的含义）
  \item 等高线用来在平面上表示三维地形——线越密坡度越陡，线越疏坡度越缓
  \item 经纬度系统为地球表面上的每个点提供了唯一坐标——经度决定时区，纬度影响气候
\end{itemize}

\subsection*{讨论题}
\begin{enumerate}
  \item 墨卡托投影让高纬度地区面积严重失真——你认为这种"不准确"的地图会在人的潜意识里造成什么影响？如果换一种投影让非洲看起来巨大而欧洲很小，人们的"世界认知"会不同吗？
  \item 除了地图投影外，你还知道有哪些"看似客观中立但实际带有视角偏见"的信息呈现方式？（比如统计图表、新闻报道、摄影镜头）
\end{enumerate}

\subsection*{动手试一试}

\begin{enumerate}
  \item 用手机地图找到你家，读出经纬度坐标。然后标出学校——在两点之间拉一条直线。这段距离在地面上是多少？
  \item 找一张秦岭地区的等高线地形图。试着判断——从北坡上去比南坡更陡还是更缓？渭河在哪一边流的？在这个坡度方向感里，你能猜出古代打仗为什么会在这里死磕吗？
\end{enumerate}

\newpage

\section{\textcolor{orange!80!black}{第二课\quad 中国：胡焕庸线——一条线分成两个中国}}

\subsection*{94\%的人住在线的那一侧}

1935年，中国地理学家\textbf{胡焕庸}在论文里画了一条线。这条线从东北的\textbf{瑷珲}（今天的黑龙江黑河）拉到西南的\textbf{腾冲}（云南西部边境小镇）。线以东，面积占全国约36\%，但住着\textbf{大约94\%的中国人}。线以西，面积占全国约64\%，人口只占约6\%。八十年过去了，城镇化、高铁、西部大开发\u2014\u2014这条比例几乎没怎么移动。\textbf{它大致没变，因为它背后站着的不是政策，是自然地理。}

\subsection*{三大阶梯——中国的地形骨架}

中国的地形可以非常直观地分成三级向下的台阶：

\textbf{第一级阶梯：青藏高原。}平均海拔4000米以上，面积占全国约四分之一。这里的氧气含量只有海平面的一半多一点（约60\%）——\textbf{高原反应}让大量人口长期定居的可能性极低。青藏高原本身挡住了从印度洋北上的潮湿气流\u2014\u2014喜马拉雅南麓是全球降水最多的地方之一，但翻过山脊之后——翻进了青藏高原的北侧——降水量骤降到几乎为零。

\textbf{第二级阶梯：高原与盆地。}海拔一千到两千米——内蒙古高原、黄土高原、云贵高原、塔里木盆地、四川盆地都在这级台阶上。黄土高原覆盖着一层厚得惊人的疏松黄土（最厚处超过300米），经过黄河几千万年往下游冲刷\u2014\u2014这就是\textbf{黄河泥沙量极高}的根本原因。

\textbf{第三级阶梯：东部的平原与丘陵。}海拔500米以下——东北平原、华北平原、长江中下游平原连续分布。平坦、河流密布、水热条件够——农业可以支撑极高的人口密度。\textbf{胡焕庸线几乎就压在第一第二级阶梯向第三级阶梯过渡的生态垂落线上。}线以西\u2014\u2014要么太高（青藏）、要么太干（内蒙古和新疆沙漠化区、黄土高原的缺水旱区）\u2014\u2014农业承载不了高密度人口\u2014\u2014商业活动随之变薄\u2014\u2014城市相互隔开。

\subsection*{秦岭-淮河：北方和南方的分界线}

中国地理上还有一条东西向的分界线\u2014\u2014秦岭-淮河一线。

这条线的北侧：年降水量一般在800毫米以下——以小麦、玉米旱地作物为主。煤炭资源富集，城市冬季集中供暖是\u2014\u2014\u201c有资格\u201d和\u201c不需要\u201d的分界线。线的南侧：年降水量一般在800毫米以上\u2014\u2014以水稻为主的水田农业区，河网密集，船是传统交通工具。\textbf{食物、房屋朝向（北方坐北朝南为纳阳光，南方避西晒兼顾通风）、方言区、乃至农历节气的微妙差别\u2014\u2014全部沿着这条线一劈两半。}

\subsection*{北京为什么在那里？}

北京坐落在华北平原的最北端，背后就是燕山山脉，往西北穿过南口就是蒙古高原的边缘。北京的位置本质上是\textbf{农耕与游牧的交接点}。冷了\u2014\u2014游牧民族南下抢粮，北京卡在山口；暖了\u2014\u2014王朝北上出征游牧部落，北京是第一兵站。这也就是为什么明朝在永乐年间把首都从南京迁到北京——\textbf{\u201c天子守国门\u201d\u2014\u2014首都要亲自压在北边的防线正后方}。

同时北京在永定河的冲积扇上——有一条稳定的淡水河保障了城市的饮水供给。再加上京杭大运河可以运南方的粮食北上至北京——明朝之前元大都的粮草基本靠运河活着\u2014\u2014北京虽然没有江南上海的港口便利，但靠的是南北向的大运河强制把经济流扭向政治心点。\textbf{北京=\u200c军事壁垒+水源+运河输送。}

\subsection*{为什么有的省富、有的省穷？}

打开一张中国各省人均GDP的地图。排在前面的几乎全贴着海岸线——广东、江苏、浙江、福建、上海。排在后面的——甘肃、贵州、云南的大片山区——与任何出海口的直线距离可能是前面那些省份的十倍以上。

广东——珠江口+南海——从1978年之后承接了第一批从香港和台湾转移过来的电子和玩具工厂。集装箱从东莞装货，开到深圳蛇口或广州南沙只需要几小时——装上远洋轮船之后直出南海进入全球航线。江苏——长江出口+上海港的腹地——苏州和无锡的工业园区紧贴着这条黄金水道。

甘肃的大部分地区被黄土高原的沟\rube{壑}{hè}和戈壁切割成运输线上的天险。从兰州发货到东部港口——不论走铁路还是公路——每一吨货物挂的运费已经把利润表的前几行刷红了。贵州——典型的喀斯特溶蚀山区——修一公里高速公路要打穿好几条隧道和高架桥。

\textbf{沿海vs内陆不是你投不投入——是物理距离决定了你每卖一件产品要多烧多少吨柴油。}

\subsection*{要点总结}
\begin{itemize}
  \item 胡焕庸线（黑河—腾冲线）以东36\%的国土住着约94\%的人口——这条线八十年几乎没变，因为它背后的决定因素是自然地理（降水、地形、气温）
  \item 中国地形分为三大阶梯：青藏高原（第一阶梯）→高原盆地（第二阶梯）→东部平原丘陵（第三阶梯）
  \item 秦岭-淮河一线是中国南北方的分界线——以北以小麦为主、有集中供暖；以南以水稻为主、无集中供暖
  \item 沿海省份经济比内陆发达的根本原因之一是地理距离——运输成本决定了内陆省份每出口一件产品要多花更多钱
\end{itemize}

\subsection*{讨论题}
\begin{enumerate}
  \item 胡焕庸线八十年几乎没变——你觉得未来有没有可能因为技术进步（比如远程办公、海水淡化、垂直农业）而改变？如果能，会是什么形式？
  \item 沿海和内陆的经济差距很大程度上是地理造成的——你觉得政府应该通过转移支付来"补"内陆，还是鼓励人口向沿海迁移？这两种思路各有什么利弊？
\end{enumerate}

\newpage

\section{\textcolor{orange!80!black}{第三课\quad 日本：没有资源的国家怎么活下去？}}

\subsection*{山脉占了国土的三分之二}

日本列岛是\textbf{太平洋板块和菲律宾海板块俯冲到亚欧板块下方}直接挤出来的火山弧。全国大约73\%的国土是山地。平原极少且狭长——最大的关东平原（东京所在地）面积也不过约1.7万平方公里——相当于北京市的面积（1.68万平方公里）略大一点点。

而这个国土上的资源极其匮乏：石油自给率在0.5\%以下，铁矿石几乎没有，煤炭在北海道和九州有一些但开采成本过高，已基本关停。\textbf{原材料全靠船运进来——产品再靠船运出去。}

\subsection*{太平洋工业带}

这一切逼出了日本经济地理上最显著的地图特征——\textbf{太平洋带状工业地带}。从东京湾的京滨工业区，经名古屋的中京工业区、大阪神户的阪神工业区，到北九州市和福冈——几乎所有重化工业全部贴在太平洋和濑户内海的港口岸线上。

逻辑极其简单：\textbf{工厂不能离港口远，因为原材料和燃料来自万吨巨轮的船舱。}铁矿石从澳大利亚和巴西用巨型矿砂船运到日本的沿海钢厂——船直接贴到工厂码头上，矿砂由传送带直接送进高炉——如果工厂在内陆，日本人需要先卸船、装卡车或内河船、再卸车一次——成本成倍翻升。

鹿儿岛县有一个极端的案例——\textbf{濮物临海钢铁厂}——它完全是在海上填出来的。人工岛直接建了高炉、轧钢车间和专用的深水泊位\u2014\u2014巨大的矿石船停靠后卸货距离不超过几百米。

\subsection*{渔场——亲潮和黑潮的交汇}

日本列岛的东北方向，从千岛群岛下来的\textbf{亲潮}（冷水流）与从南边赤道附近涌上来的\textbf{黑潮}（暖水流）在日本东北部的三陆冲海域相遇。冷暖交汇处\u2014\u2014大量的浮游生物被翻搅到水面，小鱼吃浮游生物，大鱼吃小鱼——\textbf{形成了世界三大渔场之一：西北太平洋渔场}。日本人的蛋白质结构从鱼而来——这就不是偶然的。

\subsection*{美国的"小麦带"（对照）}

横向对比一下美国和日本的农业地理，能把两边的自然约束看得更清楚。

美国本土的大平原\u2014\u2014从得克萨斯向北穿过堪萨斯、内布拉斯加到南北达科他\u2014\u2014是大规模机械农业的天然实验室。\textbf{地形平坦允许巨型联合收割机直线开行，温带大陆性气候夏季高温且有一定降水利于冬小麦和春小麦的周期种植}。美国一个农民平均经营面积可以达到上百公顷——在日本\u2014\u2014多山、田地小碎块、水稻田低洼排水不畅\u2014\u2014相同土地面积下需要几十户甚至上百户投入分散劳力。

美国的小麦带往北一些（气候更凉）是玉米带——爱荷华、伊利诺伊、印第安纳——玉米被大量处理成饲料转化为肉类和乙醇。\textbf{两种大宗商品完整支撑了美国的食品工业和出口结构——这全部基于"平坦+黑土+够用降水"这把地理牌。}

\subsection*{要点总结}
\begin{itemize}
  \item 日本73\%的国土是山地，资源极度匮乏——石油自给率不足0.5\%，铁矿石几乎为零
  \item 太平洋带状工业地带把重化工厂全部建在港口边，因为原材料全靠海运——工厂离港口越近成本越低
  \item 日本东北海域冷流（亲潮）与暖流（黑潮）交汇形成了世界三大渔场之一——这是海洋地理的馈赠
  \item 对比美国：平坦的大平原让大规模机械化农业成为可能，而日本多山碎块的地形决定了精细小农模式
  \item 韩国和越南走的是不同的"靠海生存"之路：韩国把货值压缩到芯片（运费占比极小），越南卖劳动力承接组装线
\end{itemize}

\subsection*{讨论题}
\begin{enumerate}
  \item 如果日本没有太平洋沿岸的深水良港，它的工业历史会有多不同？你觉得一个国家的"地理运气"在其发展中占多大比重？
  \item 韩国选择把资源集中在芯片这种"高附加值、轻重量"的产业上——如果当初韩国选择了钢铁和造船而不是半导体，它的今天会怎样？产业选择是"规划出来的"还是"被地理逼出来的"？
\end{enumerate}

\newpage
\subsection*{东亚邻居——韩国和越南，同一条地理剧本的不同答案}

韩国和日本的地理约束几乎是从同一副模具里倒出来的：国土多山、平原狭小、铁矿石和石油几乎全靠进口。韩国把这条剧本演成了自己的版本——\textbf{三星}从一家卖干鱼和蔬菜的小贸易商社转型为全球最大的半导体和智能手机制造商之一；\textbf{SK海力士}把存储芯片做成了能与美国和日本巨头同台竞争的产品。和日本相比，韩国政府把资源更集中地倾注给了极少数家族财阀——用政府担保的贷款和外汇政策保护它们对冲进口原材料的成本压力。半导体和智能手机不像钢铁那样按吨位计价——一块晶圆的重量不过几克，但它上面刻着几十亿个晶体管。\:\textbf{韩国的地理答案不是"在港口旁边建一个钢铁帝国"，而是"把你的货值压缩到每克能卖几百美元的芯片上——运费已经被消灭了"}。

越南自1986年"\textbf{革新开放}"之后走上了一条不同的岔路。它紧贴着中国的南大门——当制造商在中国沿海的工资成本逐步上升时，\textbf{越南用更年轻、工资更低的劳动力把一部分服装、鞋类和电子组装线吸到了自己的南海海岸线上}。越南的海岸线长达3260公里，从海防到胡志明市的货轮可以贴着海岸线一路开到新加坡海峡再接入全球航线。\textbf{日本卖技术、韩国卖芯片、越南卖劳动力——同一张"多山少资源靠海生存"的地理考卷，交出了三种完全不同的答卷。}

\newpage

\section{\textcolor{orange!80!black}{第四课\quad 美国：两洋之间的"地理彩票"}}

\subsection*{一种全球独一份的地理位置}

打开世界地图，看北美洲的那一大块叫美国的区域。它的东西两侧\u2014\u2014一边是大西洋，一边是太平洋\u2014\u2014两边都是宽阔得足以用来阻挡入侵的天险海洋。北边的邻居加拿大人口稀少且不动武，南边的墨西哥\u2014\u2014地理上被沙漠和格兰德河隔开。在十九世纪迅速扩张之后，美国实质上获得了一个\textbf{没有任何邻近强国可以直接威胁它核心领土的堡垒位置}。两次世界大战期间，美国本土没有落过一发炸弹——而在欧洲和亚洲的整个战场上，城市被夷为废墟。\textbf{不受本土打击就可以持续地生产武器、资源和粮食，这是它最深层的战略底牌。}

\subsection*{密西西比河——老天爷铺的水路网络}

密西西比河和它的主要支流（密苏里河、俄亥俄河）构成了一张\textbf{内陆水运大网}。这条河的主流加支流通航总长度约两万多公里，把中西部大平原的玉米和大豆用成本极低的驳船运到新奥尔良港口\u2014\u2014再接驳海船出口全球。水路运输的廉价程度用一组简单数字就能看清楚：同样一吨货物，运一公里的内河驳船燃油成本只相当于铁路的三分之一不到，公路的一成。\textbf{密西西比河不是美国人挖出来的——是全球几百万年的冲积送给了那片大陆。}

\subsection*{五大湖锈带——把铁、煤和水焊在一起}

美国东北部靠着加拿大边界\u2014\u2014苏必利尔湖、密歇根湖、休伦湖、伊利湖、安大略湖——五大湖构成了一整套\textbf{天然联运系统}。

苏必利尔湖的西端贴着明尼苏达的\textbf{铁矿石矿脉}（梅萨比岭），宾夕法尼亚州和西弗吉尼亚州下面压着巨量的\textbf{石炭纪产出优质炼焦煤}。煤和铁中间\u2014\u2014隔着五大湖的便利水运\u2014\u2014把两个重吨位原矿用货运轮船大量对拉。芝加哥和底特律到克利夫兰、匹兹堡和布法罗产生了闻名世界的\textbf{五大湖"锈带"}——汽车制造、钢铁、重型机械曾经在一百年里密集地锁在这一大圈水道上。\textbf{当一个地方同时凑齐了：铁矿+煤矿+运矿的水路——大规模钢铁生产就会自己长在那里。}

\subsection*{硅谷为什么不会在爱荷华？}

加州的\textbf{硅谷}（从旧金山湾往南延伸到圣何塞的那一片平坦谷地）的生产要素不是铁矿——是\textbf{斯坦福大学}的工科研究引擎、太平洋另一端的亚洲元件供应链、美国军方的早期半导体订单和\u2014\u2014\u201c不结冰的加州气候\u201d\u2014\u2014这些全部粘在了同一块地理芯片上。\textbf{高技术产业的聚集不是因为地下有煤——它聚集在资本和研究机构密集、工程师愿意搬来生活的地方。}

\subsection*{一张卫星灯光图讲完}

如果晚上从太空俯拍美国本土——你会看到明暗极其不均。东海岸从波士顿到华盛顿一条密集光斑串（人口走廊）；五大湖一圈明亮的灯斑；加州沿海一条细细的发光带；德克萨斯州-休斯顿-达拉斯呈双核闪亮。剩下的\u2014\u2014从中西部的达科他到蒙大拿、怀俄明、内华达\u2014\u2014大片的暗黑色块：人口密度极低。\textbf{卫星灯光是所有自然地理把人类安置到特定几个褶皱带的最诚实的目击者。}

\xuehai{集聚经济——为什么工厂喜欢扎堆}{
  经济学家发现企业在地理上"扎堆"有三个独立的好处：第一，共享劳动力池——硅谷的工程师在这家公司干腻了可以走两条街去另一家，公司裁员时被裁的人也容易在本地找到下一份工作。第二，共享供应商——一家汽车厂周围会自发长出轮胎、座椅、车灯和物流公司，因为谁也不敢把零件厂开在远到运输成本压死自己的地方。第三，知识溢出——底特律的工程师跳槽时把上一家研究的发动机扭矩数据装在脑子里带到了下一家的咖啡间。这三件东西——人、零件和脑浆——不会通过电子邮件远距离传递。它们会坐着同一趟地铁、穿过同一个工业园区的大门。
}

\subsection*{要点总结}
\begin{itemize}
  \item 美国拥有全球独一无二的"堡垒"地理位置——两面临洋、邻国不构成威胁，本土在两次世界大战中未被攻击
  \item 密西西比河水系构成了美国内陆运输的大动脉——水运成本仅为铁路的三分之一、公路的十分之一
  \item 五大湖区同时凑齐了铁矿、煤矿和水运——钢铁工业"自己长在了那里"
  \item 硅谷的崛起靠的不再是自然资源，而是大学（斯坦福）、风险资本、国防订单和宜人的气候
  \item 卫星灯光图忠实地反映了自然地理对人类活动的约束：人口和产业集中在少数几个"褶皱带"
\end{itemize}

\subsection*{讨论题}
\begin{enumerate}
  \item 美国的地理"安全"（两洋屏障、弱邻）在多大程度上塑造了它的国际行为方式？如果美国位于欧洲那样的位置上（强邻环伺），它的外交和军事策略会有多不同？
  \item 集聚经济（agglomeration）解释了为什么工厂喜欢扎堆——你觉得这个逻辑在互联网时代还成立吗？远程办公会不会让硅谷式的地理集聚变得不再必要？
\end{enumerate}



\newpage

\section{\textcolor{orange!80!black}{第五课\quad 西欧：河流、煤田与欧盟}}

\subsection*{莱茵河——欧洲的货流大动脉}

西欧的内陆没有一条河能和莱茵河相较。它从瑞士阿尔卑斯山的冰川融水起源，经过德国、法国边界，穿过荷兰的三角洲最终注入北海。从巴塞尔（瑞士的莱茵河航运起点）到鹿特丹港\u2014\u2014全程约800公里全年可通航\u2014\u2014将瑞士、德国、法国和荷兰用一条成本极低的水上货道贯通。

\textbf{鲁尔区}恰好位于莱茵河与其支流鲁尔河的交汇带。鲁尔区地下埋着德国最大的优质硬煤——河能从外部把瑞典的铁矿石运进来。煤+铁+一条能扛得动百万吨原料的航运动脉——和美国的五大湖铁锈区几乎是完全对称的经济地理故事。

\subsection*{鹿特丹——一个河口港如何服务整个大陆}

莱茵河最终流进荷兰的鹿特丹港。鹿特丹不是荷兰自己一国的港——它是整个西欧输入大宗产品的大门。巨型油轮和矿砂船在鹿特丹卸货，货物随即被换装上莱茵河的驳船，逆流推往杜伊斯堡、科隆、曼海姆、巴塞尔。\textbf{鹿特丹每年处理的货物以数亿吨计}——它的吞吐量常年在全球排名前列，而荷兰本国人口只有不到两千万。

\subsection*{伦敦为什么在那里？}

伦敦坐落在泰晤士河的河口——从北海进入英国内部的第一道深水港湾。泰晤士河潮汐每天上涨下降把河口拉宽\u2014\u2014让伦教可以停泊大量吃水深的大船。\textbf{面向大西洋的方位在英国殖民时代又是它指挥全球贸易的船站的天然位置}。城市本身不产煤、不产铁——但它不需要产——它只要占据收过路费和调度全球利润的位置就够了。

\subsection*{巴黎——法兰西岛的单核辐射}

法国地图最独特的一点是\textbf{极度单核化}。巴黎坐落在塞纳河的一个弯曲内——从古罗马时代就是法兰西岛区域政治控制的核心。一千多年下来所有铁路、高速公路、教育和经济决策机构全面向巴黎集中堆叠。今天你在法国地图上用人口线去量——\textbf{巴黎都会区人口超过一千二百万，占了全国总人口的近六分之一。}第二大城市里昂的人口大约百分之十出头。

单核结构的脆弱面非常明显：如果巴黎瘫痪\u2014\u2014整个国家的指挥链条断裂。

\subsection*{欧盟——把煤和铁锁在一起}

1951年，法国、德国、意大利、比利时、荷兰和卢森堡六国签署了《\textbf{欧洲煤钢共同体}》条约——把煤和钢的生产置于一个共同的超国家机构监督之下。选择煤和钢的意图完全不掩饰：\textbf{要造坦克和大炮，这两个原材料缺一不可。}如果把煤钢的生产互相绑在一起，任何一个成员国悄悄扩军的行动都会被共享信息、连于同一份账本的链条提前发现。

这个煤钢共同体直接演化成了今天的\textbf{欧盟}核心\u2014\u2014\:二十多个国家之间免关税、免护照、统一用欧元（多数成员国）。从打了两次世界大战血流成河，到互免签证去隔壁国买菜——地理在这个进程里的角色永远是：西欧的地块太小，工业带互相卡在一起——\textbf{物理密集本身逼迫合作和共享资源。}

\subsection*{要点总结}
\begin{itemize}
  \item 莱茵河从瑞士阿尔卑斯山流到北海，全年通航约800公里——是欧洲最核心的水运大动脉
  \item 鲁尔区与美国的五大湖锈区几乎是对称的经济地理故事：煤+铁+水路=重工业自然生长
  \item 鹿特丹港服务的是整个西欧内陆（而非仅荷兰一国）——这是莱茵河的入海口赋予它的功能
  \item 巴黎的极度单核化（全国近六分之一人口聚集）是千年政治集权在地理上的投射
  \item 欧洲煤钢共同体→欧盟的演化逻辑：地理逼迫西欧各国共享资源以避免互相毁灭的战争
\end{itemize}

\subsection*{讨论题}
\begin{enumerate}
  \item 巴黎的单核辐射对法国是好事还是坏事？如果你是法国决策者，你会想办法"分散"人口——还是接受这种集中并让它变得更高效？
  \item 欧洲煤钢共同体的设计逻辑是"把战争原料锁在一起，让战争变得不可能"——你认为这种思路在今天的国际关系中还能应用吗？比如在能源或芯片领域？
\end{enumerate}

\newpage

\section{\textcolor{orange!80!black}{第六课\quad 俄罗斯：最大的国家，最难的地理}}

\subsection*{全球最大的陆地拼图——但大多是冻土}

俄罗斯的面积约1709万平方公里，接近地球陆地总面积的八分之一。它11个时区\u2014\u2014当莫斯科人上午九点准备开会时，堪察加半岛的居民已经在下午六点做晚饭了。

但是翻开气候带图和冻土分布图——\textbf{西伯利亚的大部分地面是永久冻土}。夏季表面层融化但下去几十厘米之后土是千年不化的冰固层。在这种地面挖地基、铺铁路、搭油管都是工程学上极其昂贵和危险的——所有金属在反复冻胀后会变形。

\subsection*{不冻港的百年执念}

俄罗斯的海岸线极长，但\textbf{全年不结冰的温水港少得可怜}。波罗的海的圣彼得堡纬度接近60度，冬天海面结冰需要破冰船强制开路。符拉迪沃斯托克（海参崴）也需要冬季使用破冰船或等待破冰轮扫出水道。黑海的塞瓦斯托波尔和诺沃罗西斯克是为数不多的几个相对温水港——但要从黑海进入地中海，船必须穿过\textbf{博斯普鲁斯海峡和达达尼尔海峡}——两条水道全部是土耳其的领海（国际公约控制通行但战时处境会急剧变化）。这解释了俄罗斯为什么\u2014\u2014不是由于某个领袖的偶然野心\u2014\u2014持续几百年试图把势力推到克里米亚半岛的南端。

\subsection*{西伯利亚大铁路——一根横贯大陆的钢索}

从莫斯科到符拉迪沃斯托克——\textbf{西伯利亚大铁路}全长约9289公里，横贯八个时区，1891年动工，1916年完成。这条铁路不是"为了方便旅游"。在它铺完之前，俄国远东的部队和物资从欧洲部分调动需要绕行全球海运数月\u2014\u2014有了这条铁路之后，货物和军队可以\textbf{在几周内由陆路搬到亚洲海岸}。这是一条\textbf{用铁轨强行缝合了被距离撕裂的巨大领土}的生命线。

\subsection*{能源出口——经济命脉系于管道}

俄罗斯是全球最大的天然气出口国之一，也是最大的石油出口国之一。西伯利亚的气田（超大型的乌连戈伊气田）通过直径超过一米四的长输管线把天然气送往欧洲和东方。\textbf{能源出口额为联邦预算贡献了重要份额}——地理决定了地下是什么，但出口依赖也意味着当买家转向替代供应时，俄罗斯承受的冲击会格外沉重。

\subsection*{要点总结}
\begin{itemize}
  \item 俄罗斯是地球上面积最大的国家（1709万平方公里，11个时区），但大部分领土是永久冻土——工程建设和居住成本极高
  \item 不冻港的稀缺是俄罗斯几百年来持续向南扩张的核心地缘驱动力——不是某个领袖的野心，而是地理在驱动战略
  \item 西伯利亚大铁路（9289公里）是一条用铁轨强行缝合被距离撕裂的巨大领土的生命线
  \item 俄罗斯经济高度依赖能源出口——地理决定了地下的资源禀赋，但也使经济结构单一化、易受外部冲击
\end{itemize}

\subsection*{讨论题}
\begin{enumerate}
  \item 如果俄罗斯拥有像法国或英国一样温暖的不冻海岸线，它的历史会完全不同吗？你觉得"追求不冻港"在多大程度上塑造了俄罗斯的国家性格和国际行为？
  \item 一个国家的领土"大"一定是好事吗？俄罗斯的例子告诉我们——领土越大，维护成本越高，而且越大可能越脆弱。你能想到其他"大反而带来麻烦"的国家的例子吗？
\end{enumerate}

\newpage

\section{\textcolor{orange!80!black}{第七课\quad 中东：石油、沙漠与咽喉要道}}

\subsection*{石油的分布——半个世界的储量锁在这里}

根据BP世界能源统计年鉴的数据，中东地区已探明的石油储量占全球\textbf{接近一半}。最集中分布在波斯湾沿岸的巨型油田群：沙特阿拉伯的\textbf{加瓦尔油田}（全球最大陆上油田，探明可采储量达百亿桶级别），科威特的布尔甘油田，伊朗和伊拉克交界的多个超大型含油构造区。

为什么这里这么多油？地质上——\textbf{中生代到新生代的特提斯洋}（古地中海）在闭合过程中的温暖浅海环境沉积了大量有机质，加上合适的地层圈闭构造——把有机质闷在高温高压下变成了油气藏。沙特一个油田一天可以自喷上百万桶——\textbf{抽上来几乎不需要注水加压}\u2014\u2014开采成本极低。

\subsection*{霍尔木兹海峡——全球油路的"针眼"}

在波斯湾的出口处，伊朗和阿曼之间的那一道窄窄的水道——最窄处仅约\textbf{33公里}宽——是\textbf{霍尔木兹海峡}。全球每天海上贸易石油的约五分之一在这道海峡两侧互相通行。如果因为冲突被封锁\u2014\u2014国际油价会在几天之内飞升，所有依赖从中东进口原油的炼油厂从亚洲的日本、中国、印度，到欧洲的意大利和西班牙都将面临原料突然中断。

\subsection*{石油和水——沙漠里更大的问题是水}

地图上几乎整个阿拉伯半岛都埋在沙漠气候下——年降水量不足100毫米。沙特阿拉伯没有永久性河流。城市和农业用水——从首都利雅得到麦加——绝大部分来自\textbf{海水淡化工厂}和深层\textbf{地下水（化石水\u2014\u2014几千年前积攒下来的不可再生的地下水）}。一个国家同时卖全人类最贵的出口商品之一（石油）\u2004\u2014\u2004然后拿卖油的收入的一部分去淡化\u2014\u2014\textbf{淡化一桶水消耗的能源比你手上那桶原油的当量差不了多少。}

\subsection*{苏伊士运河——一条水道堵住地球}

1869年凿通的\textbf{苏伊士运河}连接了地中海和红海\u2014\u2014从伦敦到孟买的航程比绕行非洲好望角节省了约\textbf{7000公里}。全球货物贸易中大约有十分之一通过苏伊士运河。2021年3月，超大型集装箱船"\textbf{长赐号}"在运河南段搁浅\u2014\u2014\textbf{堵死了双向全部通行整整六天}。据估算，全球贸易在这六天里每一小时的堵断处理成本高达数亿美元——一条宽度只有几百米的运河南端歪了一艘船。

\xuehai{荷兰病——当一种资源"太赚钱了"会怎样}{
  1970年代荷兰在北海发现了大型天然气田，出口收入激增。但与此同时荷兰的制造业出口竞争力急剧下滑——因为天然气出口带来的大量外汇把荷兰盾（当时的货币）的汇率推高了。本币升值意味着荷兰制造的自行车和机床在国际市场上突然变贵了——买家转向了德国和日本。经济学家后来把这种现象称为"荷兰病"：一个突然暴富的资源部门通过抬高汇率，间接压垮了本国的制造业。沙特阿拉伯和俄罗斯都面对同一张压力——卖油换来的美金一旦回流本国经济、推高物价和工资，出石油之外的其他一切出口品都变贵了。\textbf{有钱不是病——但只有一条来钱的路，其他地方可能会退化成荒漠。}
}\subsection*{迪拜}——一个只剩不到1\% GDP靠石油的石油国家}

阿联酋的迪拜把一手"坐拥油田"的牌换成了一套完全不同的打法。迪拜的石油储量相对不大——政府很早就算过一笔账：\textbf{靠卖油撑不了几十年}。于是他们把卖油挣到的第一桶金塞进了几个方向：建全球最大的航空枢纽之一——\textbf{阿联酋航空}把迪拜变成了任何从欧洲飞往亚洲或大洋洲几乎绕不开的中转站；建杰贝阿里港——中东最大的人造深水港，加上紧贴港口的自由经济区，外国公司可以在这里设厂免税运营；再铺了一整条旅游和金融服务的地毯。迪拜做这件事不是因为"石油国家天然爱服务"——是管道不会说话。石油会卖完，但跑道和港口不会突然从地图上消失。\textbf{资源型国家在资源耗尽之前，可以选择把钱换成不可撤回的实体基础设施。}

\subsection*{要点总结}
\begin{itemize}
  \item 中东地区已探明石油储量占全球接近一半——地质上这是因为古特提斯洋在闭合过程中形成的有利成油条件
  \item 霍尔木兹海峡最窄处仅33公里，全球每天约五分之一的海运石油通过此处——它是全球能源供应链最脆弱的一环
  \item 沙特等海湾国家极度缺水——用水靠海水淡化和不可再生的化石地下水，这需要消耗大量能源
  \item "荷兰病"警示资源型国家：单一资源出口会推高汇率、压垮其他产业——丰富可能变成诅咒
  \item 迪拜提供了另一种可能：在石油耗尽之前把收入转化为不可撤回的实体基础设施（港口、航空枢纽、金融中心）
\end{itemize}

\subsection*{讨论题}
\begin{enumerate}
  \item 如果中东的石油明天突然枯竭——那些靠石油收入维持的国家会发生什么？你能想到哪些国家已经提前做了"后石油时代"的布局？
  \item "荷兰病"说明资源丰富不一定是好事——你能想到现代社会中有类似的"单一依赖"现象吗？（不仅限于自然资源，也可以指城市对某一产业的依赖）
\end{enumerate}

\newpage

\section{\textcolor{orange!80!black}{第八课\quad 澳大利亚与全球贸易网}}

\subsection*{一整个大陆的资源，只给几千万人用}

澳大利亚的面积约769万平方公里——差不多相当于美国的本土面积（不含阿拉斯加）——但人口只有约2600万。\textbf{大片的土地是沙漠和半干旱灌木区}——中部被称为"\textbf{内陆}"（Outback）的地区年降水量极低。人口几乎**全部挤在东南沿海的几个大城市**——悉尼、墨尔本、布里斯班——以及西南角的珀斯。从悉尼到珀斯的直线距离已超过3200公里——和一个从北京直飞乌鲁木齐差不多了——但是在同一个国家。

人口这么少却要维持一个大陆的基础设施——这是个代价极高的成本摊派。跨东西的铁路干线（印度洋-太平洋铁路线，全长约4352公里）一年运送的旅客和货物摊算分摊到每一个纳税人身上——你交了相当一笔税金就为了\textbf{保证一个两头巨型空间的联结不被拉断}。

\subsection*{铁矿、煤炭和液化天然气——"幸运之国"的出口篮子}

西澳大利亚州的\textbf{皮尔巴拉地区}铁矿石的含铁品位极其优越，普遍达到60\%以上——很多矿石不用经过富集选矿就可以直接进高炉炼钢。这些巨大的红色矿体\u2014\u2014和巴西的卡拉加斯矿山\u2014\u2014用世界上最大的散货船（好望角型、巨型矿砂船）运往中国、日本和韩国的沿海钢厂。

昆士兰州和新南威尔士州的优质动力煤和炼焦煤同样大量出口。液化天然气方面\u2014\u2014西北大陆架的天然气经过零下162度的液化后装进专门的LNG船运往东北亚——\textbf{天然气不用管道跨海}——这是澳大利亚能把自己气体的分子从南半球卖到北半球的需求区里最重要的物理条件。

\subsection*{集装箱革命——一个铁箱子是怎么把全球运费打下来的}

直到1956年以前，港口卸货这一过程是\textbf{极其笨拙的}——成百上千个大小不一的包裹在码头靠码头工人一个一个从船舱里扛出来。美国卡车司机马尔科姆·\textbf{麦克莱恩}在1956年提出了一个完全重新设计的方案：\textbf{把货物装进一个标准尺寸的铁箱——箱体从卡车拖头直接吊装到轮船专用的叠放架上}——到达目的港之后直接吊装到另一辆卡车或火车底盘。\textbf{装卸从"扛包"变成了起重机的一次垂直抬放}。一吨货物从上海漂到鹿特丹（约两万公里的远洋航线）\u2014\u2014运费分摊下来折合不到一杯咖啡的价格。这就是为什么你脚上穿的鞋可能是越南制造、手机在中国组装、奶粉标签印着澳大利亚。\textbf{集装箱把"距离"杀死了——运费压缩到了商品价格的极小百分比。}

\subsection*{全球化的两面}

集装箱网把全世界港口串成了一个协同生产的巨网。好处是：各国的工厂能专注于自己相对最擅长的制造环节。坏处同样在同一个环节：\textbf{如果一个港口里的装缷作业因为任何原因突然中断\u2014\u2014整个链下游的生产线就可能缺料停摆。}全球化将风险同步放大了\u2014\u2014经济连上同一艘集装箱货轮\u2014\u2014当船颠簸时，全船的人一起感受不稳。

\subsection*{要点总结}
\begin{itemize}
  \item 澳大利亚面积与美国本土相当但人口仅2600万，大部分国土是沙漠——人口高度集中在东南沿海几个大城市
  \item 澳大利亚的铁矿、煤炭和液化天然气大量出口——皮尔巴拉地区的铁矿石含铁量超过60\%，是世界上品质最好的矿石之一
  \item 集装箱革命把海运运费压缩到几乎可以忽略不计——这是全球化的物质前提，一吨货物从上海到鹿特丹的运费折合不到一杯咖啡的价格
  \item 全球化的两面：分工协作提高效率，但风险也同步放大——一个环节中断可能波及整个链条
  \item 巴西和澳大利亚起点相似但结果不同——说明了政治稳定性、制度可信度和与世界主航道的距离同样重要
\end{itemize}

\subsection*{讨论题}
\begin{enumerate}
  \item 集装箱让全球运费降到极低——如果能源价格大幅上涨，导致运费翻几倍，你觉得全球化会"倒退"吗？就近生产（本地制造）会不会重新变得有吸引力？
  \item 巴西和澳大利亚的对比说明：同样的自然资源不一定带来同样的发展结果。你认为"制度质量"（法治、产权保护、政治稳定）在经济发展中的作用比自然资源更大吗？为什么？
\end{enumerate}

\subsection*{同一条起跑线，不同的轨迹——巴西对照}

南半球还有另一个和澳大利亚惊人相似的国家——\textbf{巴西}。面积约851万平方公里，地广人稀，靠大宗商品出口赚取外汇：\textbf{铁矿石}（淡水河谷公司是全球最大的铁矿石生产商之一）、\textbf{大豆和牛肉}、近海新发现的盐下油田。地理起跑线几乎和澳大利亚一模一样——南半球资源大国、人口集中在东南沿海的几个城市（圣保罗和里约热内卢）。

但是巴西的人均GDP曲线和澳大利亚拉开了明显的距离。澳大利亚的政局自联邦成立以来几乎没有中断过选举循环，产权制度清晰。巴西在20世纪后半叶经历了数次政权更迭——军人独裁、货币危机——投资者不敢把钱押在明天仍然延续的制度承诺上。另外，\textbf{巴西和全球核心消费市场（北美、东亚、西欧）之间的海上距离比澳大利亚远不成比例}——从巴西桑托斯港到上海的航程，几乎是从澳大利亚皮尔巴拉港出发的一倍。

\section*{讨论提示}

\subsection*{第一课（读地图）讨论题提示}
\begin{itemize}
  \item 关于"地图投影的偏见"：可以思考——每一种地图投影都有变形。墨卡托投影放大了欧洲和北美，让它们看起来比实际"重要"。地图不是中立的——它反映制图者的视角和目的。你在生活中遇到的任何"信息呈现"（图表、统计数据、新闻报道），都可以问一句：它用了什么"投影"？
  \item 关于"看似客观的信息呈现方式"：可以思考——统计图表的坐标轴起点、饼图的颜色、新闻的标题措辞——所有这些"呈现方式"的选择都在悄悄影响你的判断。没有真正"中立"的信息呈现，只有你是否意识到了其中的选择。
\end{itemize}

\subsection*{第二课（中国地理）讨论题提示}
\begin{itemize}
  \item 关于"胡焕庸线能否被打破"：可以思考——远程办公可能让部分"数字游民"迁往西部风景好的地方，但大规模人口迁移需要经济机会。海水淡化和垂直农业在沙漠里成本太高。胡焕庸线的改变可能需要一种全新的"不依赖土地的产业模式"。
  \item 关于"补贴内陆还是鼓励迁移"：可以思考——补贴内陆意味着效率损失（在不利的地方搞生产），但保持了国土的均衡发展和社会稳定。鼓励迁移意味着更高效但可能加剧区域差距和城市病。两种思路都有道理——现实中中国采取的是"两条腿走路"。
\end{itemize}

\subsection*{第三课（日本）讨论题提示}
\begin{itemize}
  \item 关于"地理运气"：可以思考——日本没有深水良港就不可能发展出今天的工业体系。但韩国和日本"地理运气"相似，结果却不完全相同。地理是起点，不是终点——同样的地理条件下，不同的制度选择和政策方向可以走出不同的路。
  \item 关于"产业选择"：可以思考——韩国的芯片产业不是天然从地里长出来的——是政府几十年持续投入、保护本土市场和集中资源的结果。地理约束（岛国、缺资源）提供了方向，但具体路径是靠人的决策走出来的。
\end{itemize}

\subsection*{第四课（美国）讨论题提示}
\begin{itemize}
  \item 关于"地理安全与外交行为"：可以思考——美国在两次世界大战中可以"隔岸观火"再决定是否参战，欧洲国家就没有这种奢侈。这种安全使得美国倾向于"不结盟"或"选择性介入"。但如果美国的地理位置像德国一样（被强邻包围），它可能发展出完全不同的外交传统。
  \item 关于"远程办公与集聚经济"：可以思考——远程办公确实削弱了面对面的需求，但集聚经济的三个好处（劳动力池、共享供应商、知识溢出）中，"知识溢出"最难被远程替代——很多创新发生在茶水间的偶遇中。所以远程办公可能让部分工作分散，但创新中心仍然倾向于集聚。
\end{itemize}

\subsection*{第五课（西欧）讨论题提示}
\begin{itemize}
  \item 关于"巴黎的单核辐射"：可以思考——集中带来效率（资源集中使用、基础设施不分散），但加剧了区域不平等。法国尝试过" decentralization"（权力下放）但效果有限，因为市场力量天然倾向于集中。也许问题不是"要不要集中"，而是"集中的收益能否通过财政转移分享给外围"。
  \item 关于"用经济绑定防止战争"：可以思考——煤钢共同体的逻辑在能源领域仍然适用：如果欧洲能源管网高度互联，任何一国切断供应的代价都包含"自己也会受损"。但在芯片领域，因为技术保密和安全顾虑，各国不太愿意"绑定"。这个思路的适用性取决于行业的"可捆绑性"。
\end{itemize}

\subsection*{第六课（俄罗斯）讨论题提示}
\begin{itemize}
  \item 关于"不冻港与俄罗斯性格"：可以思考——俄罗斯追求不冻港可以被理解为"一个身处寒冷内陆的大国对出海口永恒的战略焦虑"。这种焦虑驱动了几百年的扩张——从彼得大帝建立圣彼得堡到克里米亚的争夺。理解地理困境是理解俄罗斯外交行为的一把钥匙。
  \item 关于"大领土是福是祸"：可以思考——加拿大、澳大利亚、俄罗斯都面临"大而疏"的问题：维护成本高、人口稀疏、内部市场整合困难。相比之下，小而紧凑的国家（如荷兰、日本）虽然资源少，但治理成本低、内部市场统一。大和小各有代价。
\end{itemize}

\subsection*{第七课（中东）讨论题提示}
\begin{itemize}
  \item 关于"石油枯竭之后"：可以思考——迪拜已经转型为服务经济，但沙特仍然高度依赖石油。转型需要时间、资本和政治意愿。沙特提出的"2030愿景"试图发展旅游和科技——但一个习惯了石油收入的社会的转型，比建几个新产业更难。
  \item 关于"单一依赖"：可以思考——不只资源型国家有这个问题。一个城市如果90\%的就业来自一家汽车厂——当这家厂搬迁时，整个城市就陷入衰退。这种"单一结构"的风险，在哪里都存在。
\end{itemize}

\subsection*{第八课（澳大利亚与全球贸易）讨论题提示}
\begin{itemize}
  \item 关于"运费上涨与全球化倒退"：可以思考——运费上涨确实会让部分"全球采购"变得不划算，但完全退回本地生产也不太可能，因为很多产品的供应链已经深度嵌入。更可能出现的不是"去全球化"，而是"区域化"——在邻近国家之间建立供应链。
  \item 关于"制度 vs 资源"：可以思考——巴西和澳大利亚的对比非常有力地说明：自然资源只是"原材料"，真正决定一个国家能否把资源变成全民福祉的是制度质量。产权保护、法治和政治稳定，决定了资源收益是落入全民口袋还是被少数人攫取。
\end{itemize}

\vspace{0.5cm}
\begin{center}
\textcolor{blue!70!black}{\large —— 八课走完地球的基本经济地理拼图。从等高线到集装箱——地图在你的手里变成了一副立体的经济开关板。——}
\end{center}

\end{document}
